\chapter{Aquisição de conhecimento}
\label{cap:aquisicao}

Num sistema pericial, a qualidade das decisões depende, maioritariamente, da
forma como o conhecimento do domínio é identificado, organizado e
representado. No presente projeto, a construção da base de conhecimento
acompanhou a própria evolução do sistema, ou seja, as regras foram definidas a
partir da identificação de situações relevantes na condução, dos dados
necessários para as analisar e das possíveis ações associadas a cada
situação.

À medida que o trabalho foi progredindo, algumas das regras definidas
inicialmente foram adaptadas com base na pesquisa realizada, na análise de
cenários e no conhecimento partilhado pelo perito. Neste sentido, a base de
conhecimento é uma estrutura em evolução, que pode ser atualizada sempre que
sejam identificadas novas condições, exceções ou relações relevantes para o
processo de decisão.

\section{Fontes de conhecimento}
\label{sec:fontes}

A primeira fonte considerada corresponde à pesquisa bibliográfica e técnica,
utilizada para aprofundar conceitos relacionados com os sistemas periciais,
sensores, segurança rodoviária e a avaliação de situações de risco. Esta
pesquisa permitiu também compreender melhor determinados indicadores
relevantes para o sistema, como o TTC, utilizado
para relacionar a distância a um objeto com a velocidade de aproximação.

A segunda fonte diz respeito ao conhecimento desenvolvido pelo grupo,
nomeadamente à identificação de elementos considerados importantes para a
tomada de decisão, às possíveis relações entre estes elementos e às ações que
o sistema poderia propor. Este conhecimento foi utilizado para construir uma
primeira versão das regras a aplicar e, simultaneamente, permitiu identificar
dúvidas a esclarecer com o perito.

Por último, tivemos o contributo do perito, não só para adquirir novas
informações e conhecimento, mas também para rever e questionar algumas
decisões inicialmente definidas pelo grupo. Esta fonte de conhecimento
permitiu esclarecer aspetos do domínio relacionados, por exemplo, com
critérios de criticidade, prioridades e número de leituras a considerar.
Desta forma, a base de conhecimento é o resultado da combinação entre
pesquisa, conhecimento desenvolvido pelo grupo e conhecimento especializado.

\section{O perito e a partilha de conhecimento}
\label{sec:perito}

Para apoiar o trabalho desenvolvido, foram realizadas reuniões com o
Professor Doutor Agostinho Gil Teixeira Lopes, Professor Auxiliar do
Departamento de Engenharia Informática (DEI) do ISEP, doutorado na área de
robótica inteligente. A primeira reunião, realizada no dia 1 de outubro de
2026, teve como principal objetivo apresentar o problema em desenvolvimento,
esclarecer algumas das dificuldades identificadas pelo grupo e reter
informação especializada sobre o domínio. Um dos aspetos discutidos esteve
relacionado com a representação do ambiente, no qual o veículo circula. Até ao
momento da reunião, uma parte significativa do trabalho estava organizada
numa árvore com diferentes caminhos possíveis, o que dificultava a definição e
a articulação das situações a reproduzir nos possíveis cenários. Durante a
reunião foi sugerida, como alternativa, a possibilidade de representar o
percurso através de uma matriz, na qual o veículo e os diferentes elementos do
ambiente pudessem ser posicionados nas vias de circulação.

Nesta nova perspetiva, as posições da matriz permitiriam representar, de
forma simplificada, diferentes elementos do ambiente, ou seja, o objeto
presente na estrada seria colocado numa determinada posição, enquanto os
sinais de trânsito, por exemplo, poderiam ser associados a uma coluna ou
localização específica, ao longo do percurso. Deste modo, a matriz
funcionaria como uma representação simplificada do cenário onde irá ocorrer a
simulação, através da qual se estabelece uma relação da posição do veículo
com os elementos que têm influência na tomada de decisão. Esta primeira
reunião permitiu também reforçar a importância das variáveis como a
velocidade e a distância, relativamente aos objetos, principalmente em
situações relacionadas com o risco de colisão. Deste ponto de vista, um dos
aspetos a ter em atenção no desenvolvimento do trabalho é perceber de que
forma a evolução da relação entre o veículo e o objeto pode alterar o nível de
risco e, consequentemente, influenciar a decisão proposta pelo sistema.

Outro ponto abordado esteve relacionado com a definição da informação
utilizada durante o processo de inferência, tendo em conta que nem todos os
elementos existentes no ambiente apresentam a mesma relevância para uma
determinada decisão. Para além desta chamada de atenção por parte do perito,
foi ainda sugerida a análise da Teoria dos Jogos, como possível abordagem para
determinadas situações de interação. Esta possibilidade foi considerada
particularmente interessante para cenários em que a decisão de um veículo
pode depender do comportamento de outros intervenientes. No entanto, nesta
fase, a Teoria dos Jogos é considerada apenas como uma possível abordagem
complementar para desenvolvimento futuro. De uma forma geral, o contributo do
perito focou-se nas seguintes questões:
\begin{itemize}
  \item o comportamento do veículo em caso de falha;
  \item a relevância dos dados de cada sensor para a tomada de decisão;
  \item critérios utilizados para distinguir a criticidade do cenário e
    ambiente;
  \item a relação entre velocidade, distância e capacidade de travagem;
  \item o número de leituras a considerar antes de tomar uma decisão;
  \item os níveis de prioridade e os possíveis critérios de desempate;
  \item o tempo de permanência dos factos na memória de trabalho.
\end{itemize}

As questões colocadas ao perito foram preparadas com o objetivo de adquirir
conhecimento que pudesse, posteriormente, ser utilizado na representação e
nas regras a utilizar pelo sistema. Neste sentido, procurou-se não só
confirmar se determinadas decisões eram adequadas, mas também perceber quais
as variáveis que deviam ser avaliadas, que condições poderiam alterar uma
decisão, que exceções deveriam ser previstas e de que forma situações de
conflito ou maior risco deveriam ser tratadas.

Assim, a aquisição de conhecimento junto do perito decorreu em duas sessões.
A primeira (S1) foi uma reunião via Teams (\cref{tab:sessao1}) e permitiu
apresentar o projeto,
discutir a representação do ambiente, os factos considerados, os cenários
previstos e algumas opções de implementação. A segunda (S2), para facilitar a
comunicação com o perito, foi um questionário escrito (\cref{tab:sessao2}),
enviado por
\emph{email}, com perguntas preparadas pela equipa a partir das questões que
tinham ficado em aberto durante a construção da base de conhecimento.

\begin{table}[H]
  \centering
  \caption{Sessão 1 com o perito}
  \label{tab:sessao1}
  \small
  \begin{tabularx}{\textwidth}{@{}L{3.2cm}VX@{}}
    \toprule
    \multicolumn{2}{@{}l}{\textbf{Sessão 1 via Teams -- Reunião inicial}} \\
    \midrule
    \textbf{Data e hora} & 1 de outubro de 2026 -- 20h10 \\ \linhasep
    \textbf{Objetivo} & Apresentar o projeto ao perito e discutir a respetiva
      implementação. \\ \linhasep
    \textbf{Aspetos abordados} & Foram discutidos aspetos relacionados com a
      estrutura geral do sistema, a representação do ambiente de condução, a
      definição dos factos da base de conhecimento e a influência desses
      factos no processo de inferência e na tomada de decisão. \\
    \bottomrule
  \end{tabularx}
\end{table}

\begin{table}[H]
  \centering
  \caption{Sessão 2 com o perito}
  \label{tab:sessao2}
  \small
  \begin{tabularx}{\textwidth}{@{}L{3.2cm}VX@{}}
    \toprule
    \multicolumn{2}{@{}l}{\textbf{Sessão 2 via \emph{email} -- Esclarecimento de dúvidas}} \\
    \midrule
    \textbf{Data} & 2 de outubro de 2026 \\ \linhasep
    \textbf{Objetivo} & Esclarecer as questões em aberto sobre falhas,
      critérios de criticidade, confirmação de leituras, prioridades, memória
      de trabalho, dados contraditórios, cenários e explicações. \\ \linhasep
    \textbf{Aspetos abordados} & Neste segundo momento, foram abordados
      aspetos relacionados com a falha do veículo, relevância dos sensores,
      definição de situações críticas através de variáveis como velocidade,
      distância e TTC, número de leituras necessárias para a tomada de
      decisão, prioridades entre ações, gestão da memória de trabalho,
      tratamento de dados contraditórios, definição dos cenários de simulação
      e mecanismos de explicação \emph{Why} e \emph{Why Not}. As perguntas e
      respostas encontram-se no Anexo~\ref{anx:sessoes}. \\
    \bottomrule
  \end{tabularx}
\end{table}

O conhecimento obtido durante estas sessões contribuiu para a base de suporte
da estrutura do sistema. Alguns dos aspetos discutidos encontram-se já
refletidos no trabalho desenvolvido, nomeadamente a utilização da informação
temporal, a consideração de várias leituras em determinadas situações, a
gestão da informação presente na memória de trabalho, as prioridades nas
propostas de decisão e as explicações associadas às regras consideradas pelo
motor de inferência. Sempre que possível, procurou-se manter uma relação entre
o conhecimento adquirido, a respetiva representação através de factos e
regras e os cenários utilizados para testar o comportamento resultante. Neste
sentido, o processo de aquisição e evolução do conhecimento pode ser
representado pela \cref{fig:aquisicao}, da seguinte forma:

\begin{figure}[H]
  \centering
  \begin{tikzpicture}[
      node distance=6mm and 8mm,
      etapa/.style={draw=corEscura, fill=corFundo, thick, rounded corners=2pt,
        text width=5cm, minimum height=0.9cm, align=center, font=\small},
      fonte/.style={etapa, text width=3.6cm},
      seta/.style={-latex, thick, draw=corEscura}]
    \node[etapa] (aquisicao) {Aquisição de conhecimento\\do domínio};
    \node[fonte, left=of aquisicao, yshift=7mm] (literatura)
      {Literatura académica\\e técnica};
    \node[fonte, left=of aquisicao, yshift=-7mm] (perito)
      {Perito do domínio\\Prof.\ Gil Lopes};
    \node[etapa, below=of aquisicao] (modelo)
      {Modelo conceptual: factos,\\eventos e prioridades};
    \node[etapa, below=of modelo] (regras) {Formalização das regras};
    \node[etapa, below=of regras] (validacao) {Validação de cenários e testes};
    \node[fonte, right=of validacao] (revisao)
      {Revisão da base\\de conhecimento};
    \draw[seta] (literatura.east) -- ([yshift=2mm]aquisicao.west);
    \draw[seta] (perito.east) -- ([yshift=-2mm]aquisicao.west);
    \draw[seta] (aquisicao) -- (modelo);
    \draw[seta] (modelo) -- (regras);
    \draw[seta] (regras) -- (validacao);
    \draw[seta] (validacao) -- (revisao);
    \draw[seta] (revisao.north) |- node[pos=0.25, right, font=\footnotesize,
      text=corCinza] {iteração} (aquisicao.east);
  \end{tikzpicture}
  \caption{Processo de aquisição, validação e atualização da base de
    conhecimento}
  \label{fig:aquisicao}
\end{figure}

% ---------------------------------------------------------------
% Secção 3.3 (proposta): desativada até ser validada pela equipa.
% Para a reativar, apagar o \iffalse abaixo e o \fi no fim do ficheiro.
% ---------------------------------------------------------------
\iffalse
\section{Síntese do conhecimento adquirido}
\label{sec:alteracoes-bc}

A \cref{tab:sintese-perito} resume as indicações do perito e a forma como se
refletem no sistema, distinguindo o que já está implementado no protótipo do
que está apenas previsto. As perguntas e
respostas completas encontram-se no Anexo~\ref{anx:sessoes}.

\begin{table}[H]
  \centering
  \caption{Indicações do perito e reflexo no sistema}
  \label{tab:sintese-perito}
  \footnotesize
  \renewcommand{\arraystretch}{1.15}
  \begin{tabularx}{\textwidth}{@{}L{2.3cm}V
      >{\raggedright\arraybackslash}XV
      >{\raggedright\arraybackslash}XV
      L{2.05cm}@{}}
    \toprule
    \textbf{Tema} & \textbf{Indicação do perito} & \textbf{Reflexo no sistema}
      & \textbf{Estado} \\
    \midrule
    Falha do veículo
      & Paragem imediata (STOP)
      & Paragem com a prioridade máxima perante uma falha; o protótipo ainda
        não deteta falhas
      & Previsto \\ \linhasep
    Relevância dos sensores
      & Os dados considerados são suficientes; cada sensor é necessário na
        situação para a qual contribui
      & Sensores descritos no \cref{cap:sensores}; cinco são simulados no
        protótipo
      & Parcial \\ \linhasep
    Criticidade
      & Equação com a velocidade, a distância e a capacidade de travagem, sem
        travagens instantâneas e com um ponto de não-retorno
      & Comparação da distância ao objeto com a distância de travagem; o
        protótipo usa ainda limiares fixos de TTC e de distância
      & Previsto \\ \linhasep
    Número de leituras
      & Uma leitura numa situação crítica; duas numa situação normal
      & Uma leitura com TTC inferior a 1,3~s; duas leituras nas regras de TTC
        persistente (\cref{sec:raciocinio-temporal})
      & Implementado \\ \linhasep
    Prioridades e desempate
      & Quatro níveis (baixo, médio, alto e crítico); desempate aleatório ou
        pela primeira opção de uma lista ordenada
      & Quatro níveis de risco e sete categorias de prioridade; desempate pela
        menor velocidade recomendada (ambos justificados na
        \cref{sec:prioridades})
      & Implementado \\ \linhasep
    Memória de trabalho
      & Esquecer os eventos que saem do campo de deteção; os sinais alteram um
        estado
      & As perceções deixam de ser consideradas sem atualização durante cerca
        de 600~ms; os sinais atualizam estados, como o limite de velocidade
      & Implementado \\ \linhasep
    Dados contraditórios
      & Sem terceira fonte, aplica-se o desempate; com risco alto ou crítico,
        paragem
      & LiDAR como terceira fonte e regra de desempate
      & Previsto \\ \linhasep
    Cenários
      & Os cenários da S1, a passadeira com peão e objetos inseridos pelo
        utilizador
      & Catorze cenários e inserção de elementos pelo utilizador; a passadeira
        ainda não existe
      & Parcial \\ \linhasep
    Explicações
      & As explicações devem seguir as regras
      & Cada decisão guarda a regra de origem, a razão e a evidência
      & Implementado \\ \linhasep
    Representação do ambiente (S1)
      & Representar o percurso através de uma matriz
      & Elementos do ambiente posicionados por via e por posição ao longo do
        percurso
      & Implementado \\ \linhasep
    Teoria dos Jogos (S1)
      & Abordagem para situações de interação entre veículos
      & Considerada para trabalho futuro
      & Previsto \\
    \bottomrule
  \end{tabularx}
\end{table}
\fi
